\documentclass[11pt,a4paper]{article}
\usepackage[T1]{fontenc}
\usepackage[utf8]{inputenc}
\usepackage{lmodern}
\usepackage[margin=26mm,headheight=15pt]{geometry}
\usepackage{amsmath,amssymb,amsthm,mathtools}
\usepackage{microtype,booktabs,tabularx,longtable,array}
\usepackage{xcolor,enumitem,fancyhdr,titlesec,xurl}
\usepackage[colorlinks=true,linkcolor=navy,citecolor=navy,urlcolor=navy,
  bookmarksnumbered=true,pdfencoding=auto,psdextra]{hyperref}
\usepackage{bookmark}
\definecolor{navy}{RGB}{27,55,78}
\definecolor{muted}{RGB}{80,91,99}
\definecolor{pale}{RGB}{240,244,247}
\hypersetup{pdftitle={Sparse Computable Erasures and Least-Degree Obstructions},
 pdfauthor={Research draft prepared for Vladimir Reshetnikov},
 pdfsubject={Effective-dense reducibility, low generic sets, and Turing-degree spectra}}
\urlstyle{same}
\titleformat{\section}{\large\bfseries\color{navy}}{\thesection}{0.7em}{}
\titleformat{\subsection}{\normalsize\bfseries\color{navy}}{\thesubsection}{0.65em}{}
\setlength{\parskip}{3pt}
\setlength{\parindent}{1.2em}
\setlength{\emergencystretch}{2em}
\setlist{itemsep=3pt,topsep=5pt,leftmargin=2em}
\setcounter{tocdepth}{2}
\pagestyle{fancy}\fancyhf{}
\fancyhead[L]{\small\color{muted}Sparse erasures and effective-dense degrees}
\fancyhead[R]{\small\color{muted}Research draft}
\fancyfoot[C]{\thepage}
\renewcommand{\headrulewidth}{0.3pt}
\newcommand{\N}{\mathbb N}
\newcommand{\Cantor}{2^{\N}}
\newcommand{\Baire}{\N^{\N}}
\newcommand{\leT}{\leq_{\mathrm T}}
\newcommand{\nleT}{\nleq_{\mathrm T}}
\newcommand{\eqT}{\equiv_{\mathrm T}}
\newcommand{\ltT}{<_{\mathrm T}}
\newcommand{\degT}{\operatorname{deg}_{\mathrm T}}
\newcommand{\ned}{\mathrm{ned}}
\newcommand{\ued}{\mathrm{ued}}
\newcommand{\leN}{\leq_{\ned}}
\newcommand{\leU}{\leq_{\ued}}
\newcommand{\eqN}{\equiv_{\ned}}
\newcommand{\eqU}{\equiv_{\ued}}
\newcommand{\ED}{\operatorname{ED}}
\newcommand{\Znull}{\mathcal Z_0}
\newcommand{\Comp}{\operatorname{Comp}}
\newcommand{\Fin}{\mathrm{Fin}}
\newcommand{\Low}{\operatorname{Low}}
\newcommand{\Spec}{\operatorname{Spec}_{\mathrm T}}
\newcommand{\substar}{\subseteq^{*}}
\newcommand{\symdiff}{\mathbin{\triangle}}
\newcommand{\restr}{\mathbin{\upharpoonright}}
\newcommand{\meetstr}{\preccurlyeq}
\newcommand{\extends}{\succeq}
\newcommand{\prefix}{\preccurlyeq}
\newcommand{\pred}{\operatorname{Pred}}
\newcommand{\jumpzero}{\emptyset'}
\newcommand{\doublezero}{\emptyset''}
\newcommand{\CA}{\mathsf{CA}}
\newcommand{\DNR}{\mathsf{DNR}}
\newcommand{\nn}[1]{\lvert#1\rvert}
\newcommand{\repo}{\textsc{ProveIt}}
\newtheorem{theorem}{Theorem}[section]
\newtheorem{lemma}[theorem]{Lemma}
\newtheorem{proposition}[theorem]{Proposition}
\newtheorem{corollary}[theorem]{Corollary}
\theoremstyle{definition}
\newtheorem{definition}[theorem]{Definition}
\newtheorem{question}[theorem]{Question}
\theoremstyle{remark}
\newtheorem{remark}[theorem]{Remark}
\newtheorem{example}[theorem]{Example}
\DeclareMathOperator{\dom}{dom}
% [write] marks text added when the manuscript became report 14 (batch 37)
\newcommand{\wtag}{\textup{[write]}}
\begin{document}
% [write] the title page resets the page counter; suppress its page anchor
\hypersetup{pageanchor=false}
\begin{titlepage}
\vspace*{10mm}
{\small\sffamily\bfseries\color{navy}COMPUTABILITY THEORY\quad /\quad PROOF-FOCUSED RESEARCH\par}
\vspace{13mm}
{\Huge\bfseries\color{navy}Sparse Computable Erasures\par}
\vspace{5mm}
{\LARGE and Least-Degree Obstructions\par}
\vspace{9mm}
{\large Low counterexamples, finite meets, and descending chains\par
in effective-dense equivalence classes\par}
\vspace{12mm}
\noindent\rule{\textwidth}{0.5pt}
\vspace{4mm}
\noindent Research draft prepared for Vladimir Reshetnikov\\[3pt]
28 September 2026
\vspace{8mm}
\begin{abstract}
We give a negative answer to the effective-dense least-representative question
C2 in the \repo\ research programme, in both its uniform and nonuniform
versions. If a binary set $G$ is $1$-generic and computes no function eventually
different from every computable function, neither equivalence class of $G$
has an element of least ordinary Turing degree. Every $2$-generic satisfies
the hypothesis. A classical highness characterization also supplies all
non-high $1$-generics, including low $\Delta^0_2$ examples.

The key step turns a density-zero omission set into a bounded choice function.
An infinitely often computable agreement yields a decidable sparse mask that
escapes that omission set. Genericity then prevents recovery of an
effective-dense description from the erased oracle. For every unbounded
computable order $b$, the mask can satisfy
$|T\cap[0,N)|\leq b(N)$ at every prefix. We also identify the exact lattice
of degrees arising from computable masks, prove a uniform countable
simultaneous-avoidance theorem, and construct a strictly descending chain
inside one uniform class, all within the same budget, with no lower bound
belonging to the full nonuniform class. Every member of the chain can be low.

The main erasure argument and the $2$-generic route are proved directly.
The low-example refinement explicitly uses the classical
Kjos-Hanssen--Merkle--Stephan characterization, rather than presenting it as
new. The article records its source audit, formalization obligations, and
further questions.
\end{abstract}
\vfill
\noindent\small\color{muted}
Status: complete conventional arguments are supplied for the stated results.
This draft has not been independently refereed or checked in Lean.
The literature audit is targeted, not a certification of priority.
Finite tests in the companion archive do not certify the infinite theorems.
\end{titlepage}
\hypersetup{pageanchor=true}
\tableofcontents

\medskip
{\small\noindent
\wtag\ \textbf{Research report~14 of the CoarseDegrees programme} (batch~37, manuscript~11). Text marked \wtag{} was added when it was filed; Sections~\ref{cr14:sec:wnotation}--\ref{cr14:sec:wtargets} give the notation against the plan and the other round-2 reports, the provenance, the formal status and the plan targets addressed. The delivered \path{REPOSITORY_UPDATE.md} is filed verbatim and has not been applied. Nothing in this report is formalized. It is AI-assisted and unrefereed.\par}
\newpage

\section{The question and the results}
\label{cr14:sec:overview}
\subsection{A specific target, not a repetition of the coarse case}
The \repo\ research plan distinguishes two problems that are easy to
confuse. Its earlier coarse least-representative question, C1, is retired;
its effective-dense counterpart, C2, is retained with a proposed genericity
argument and an explicitly identified gap. The present paper addresses C2.
The repository snapshot used here is commit
\texttt{1a1396d4d3a2ac6812692df520517d86ba3a4785}; the relevant path and source
locations are recorded in Appendix~\ref{cr14:app:sources}.\cite{ProveIt}

An effective-dense description may withhold an answer on a density-zero set,
but may never give a wrong numerical answer. Withholding is represented by
an explicit symbol $\bot$, not by a computation that fails to halt. Thus
arbitrary density-zero modifications do not automatically preserve the
associated description problem. Our modifications have a \emph{computable}
support, which makes two correct, uniform description transformations possible.

For $r\in\{\ued,\ned\}$, let $[f]_r$ denote the $r$-equivalence class of
$f$, and put
\[
 \Spec([f]_r)=\{\degT(g):g\in[f]_r\}.
\]
The target is the assertion
\begin{equation}\label{cr14:eq:leastquestion}
 (\forall f\in\Baire)(\exists g\in[f]_r)
       (\forall h\in[f]_r)\;g\leT h.
\end{equation}
This is the effective-dense instance of Gerdes's Question~7. His question
also asks more generally about the spectra of ordinary Turing degrees of
representatives.\cite[Question~7]{Gerdes}
The word ``ordinary'' matters: members of $[f]_r$ have one and the same
asymptotic degree, but can have quite different Turing degrees.

\subsection{A hypothesis separating prediction from computational strength}
Say that an oracle $G$ has the \emph{computable-agreement property} if
\begin{equation}\label{cr14:eq:CA}
 \CA(G):\qquad
 (\forall f\in\Baire)\bigl[f\leT G\ \Longrightarrow\
    (\exists h\text{ total computable})(\exists^\infty n)\ f(n)=h(n)\bigr].
\end{equation}
This is simply the assertion that $G$ computes no function eventually
different from every total computable function. The notation is local to
this paper; the property itself is classical. Its highness-theoretic
interpretation appears in Section~\ref{cr14:sec:agreement}.\cite{KMS,BBNN}

\begin{theorem}[Least-degree obstruction]\label{cr14:thm:mainintro}
Let $G\in\Cantor$ be $1$-generic and satisfy $\CA(G)$.
Neither $\Spec([G]_{\ued})$ nor $\Spec([G]_{\ned})$ has a least element.
The assertion remains true when the representatives are restricted to binary
sets. Such $G$ can be chosen low and computable from $\jumpzero$.
\end{theorem}

Here ``low'' means $G'\eqT\jumpzero$. The hypotheses hold for every
$2$-generic and, by a classical theorem cited explicitly below, every
non-high $1$-generic. No assertion for \emph{all} $1$-generics is needed to
refute~\eqref{cr14:eq:leastquestion}.

For a computable set $T$, write
\begin{equation}\label{cr14:eq:maskintro}
 G_T(n)=\begin{cases}0,&n\in T,\\G(n),&n\notin T.\end{cases}
\end{equation}
The mask $T$ is a set of coordinates whose information is discarded.
The output $G_T$ is still binary. If $T$ has density zero, then
$G_T\eqU G$.

\noindent\wtag\ \emph{Not the plan's $G_T$.} The research plan's line of
attack on C2 (\path{turing_degrees_unified.tex}, line 570) writes $G_T$ for
the function equal to $2$ on $T$ and to $G$ elsewhere. Here $G_T$ is $0$ on
$T$ and stays binary. For computable null $T$ both are uniformly
effective-dense equivalent to $G$ and Turing equivalent to
$G\restr(\N\setminus T)$, so every degree statement below holds for either
reading; but the two are different functions, and only the binary one is a
set representative. See Section~\ref{cr14:sec:wnotation}.

\begin{theorem}[Arbitrarily sparse local avoidance]\label{cr14:thm:localintro}
Let $G$ satisfy the hypotheses of Theorem~\ref{cr14:thm:mainintro}.
Suppose $A\leT G$ computes an effective-dense description of $G$, and let
$b:\N\to\N$ be computable, nondecreasing, and unbounded. There is an
infinite computable density-zero set $T$ such that
\[
 |T\cap[0,N)|\leq b(N)\quad(N\in\N),\qquad
 G_T\eqU G,\qquad G_T\ltT G,\qquad A\nleT G_T.
\]
In particular, $|(G\symdiff G_T)\cap[0,N)|\leq b(N)$ for every $N$.
\end{theorem}

The candidate $A$ need not itself be an equivalent representative: computing
one effective-dense description is enough. The theorem is an existence
statement for a computable mask, not a uniform algorithm which extracts its
index from arbitrary oracle data.

\begin{theorem}[Packed descending chains]\label{cr14:thm:chainintro}
For the same $G$ and $b$, there are computable density-zero masks
$C_0\subseteq C_1\subseteq\cdots$ such that, with $G_i=G_{C_i}$,
\[
 G_0>_{\mathrm T}G_1>_{\mathrm T}G_2>_{\mathrm T}\cdots,
 \qquad G_i\eqU G,
 \qquad |(G\symdiff G_i)\cap[0,N)|\leq b(N).
\]
There is no $g\in[G]_{\ned}$ with $g\leT G_i$ for every $i$.
If $G$ is low, every $G_i$ is low.
\end{theorem}

A countable descending chain need not be coinitial in the entire class.
Nor does the last assertion deny the existence of lower bounds outside that
class: computable sets are always ordinary Turing lower bounds. Both
qualifications are essential.

\subsection{What is contributed and what is classical}
The contribution developed here is the passage from computable agreement
and generic nonprediction to least-degree obstructions for effective-dense
representatives, together with its quantitative and countable refinements.
The no-eventually-different-function fact about $2$-generics and the
high-or-DNR characterization are classical.\cite{BBNN,KMS}
We include a direct proof of the former and invoke the latter only for the
low refinement. The finite-coordinate arguments and elementary density
estimates are proved rather than treated as new background theorems.

The audit located the explicit least-representative question in Gerdes's
2025 preprint and the distinct quasiminimality statements in
Astor--Hirschfeldt--Jockusch. Quasiminimality concerns a particular embedding
of ordinary Turing degrees into asymptotic degrees; it is not the same
statement as nonexistence of a least ordinary degree \emph{within} a fixed
class.\cite[Theorem~5.9]{AHJ}
No claim of first publication or exhaustive literature coverage is made.

\subsection{\wtag\ Notation against the research plan and the other round-2 reports}
\label{cr14:sec:wnotation}
The effective-dense notions here are those of the research plan's target C2
\cite{ProveIt}: $\equiv_{\ned}$ and $\equiv_{\ued}$, with omissions written
as an explicit $\bot$ (the strong-domain convention of \cite{AHJ}). The
density-zero sets $\Znull$ are the sets the Lean library of the project
calls \texttt{DensityZero}, and ``$1$-generic'' is its predicate
\texttt{OneGeneric}. No symbol was renamed and no normalization changed
when this report was written. The other round-2 reports (11--13) are about
\emph{coarse} descriptions, which may err on a null set; several of their
letters mean something else here.

\begin{center}\small
\begin{tabularx}{\textwidth}{@{}>{\raggedright\arraybackslash}p{.14\textwidth}>{\raggedright\arraybackslash}X>{\raggedright\arraybackslash}X@{}}
\toprule
Symbol & Meaning here & Do not read it as\\\midrule
$G_T$ & $0$ on $T$, $G$ elsewhere; binary & the plan's $G_T$ (plan, line 570), which is $2$ on $T$\\
$\Spec([f]_r)$ & Turing spectrum of an \emph{effective-dense} class, $r\in\{\ued,\ned\}$ & the coarse spectra of the synthesis and of reports~11 and~13\\
$T$, $C_i$ & computable null \emph{masks}; $C_i$ is the $i$th mask of the chain \eqref{cr14:eq:chainmask} & the dyadic columns $C_i$ of the synthesis and report~13, or report~13's tails $T_k$\\
$I_n$, $J_n$ & blocks $[m_n,2m_n)$ of positions; dyadic blocks $[2^n,2^{n+1})$ in Lemma~\ref{cr14:lem:Pideal} & a block \emph{code} such as the plan's $J$, or report~13's $I_n=[2^n-1,2^{n+1}-1)$\\
$P_T$ & the projection $X\mapsto X_T$ & report~13's protected regions $P_k$, or report~11's presentation $P$\\
$r_k(n)$ & block fractions in Lemma~\ref{cr14:lem:Pideal} & report~11's column map $r_k(t)$\\
$\CA(G)$ & $G$ computes no function eventually different from all computable ones; local notation & a notion defined elsewhere in the programme\\
``low'' & $G'\eqT\emptyset'$ & report~13's ``low over $\emptyset'$'', i.e.\ $(Z\oplus\emptyset')'\eqT\emptyset''$\\
\bottomrule
\end{tabularx}
\end{center}

\subsection{\wtag\ Provenance of this report}
\label{cr14:sec:wprov}
The report is built from one manuscript: batch~37 manuscript~11, the
archive \path{proveit_effective_dense_research} (inner directory
\path{proveit_effective_dense/}, main file \texttt{article.tex}, 22-page
PDF), which arrived in commit \texttt{37e61c1fd}. Its pin is the commit
printed above, \texttt{1a1396d4d3a2ac6812692df520517d86ba3a4785} (the
batch-36 placement commit), and it was placed as report~14 in commit
\texttt{0e53d1034}. Between the pin and the placement,
\path{Computability/TuringDegrees/} changed only by the addition of the
round-2 reports 11--14, so the plan blob and the line range quoted in
Appendix~\ref{cr14:app:sources} are those of the current tree.

When the report was written, every label received the prefix
\texttt{cr14:}, the title page received the page-anchor guard that removes
a duplicate PDF destination of the delivered source, and the note below the
contents, these four subsections, and \wtag{} notes after
\eqref{cr14:eq:maskintro}, in Section~\ref{cr14:sec:gap} and in
Section~\ref{cr14:sec:checks} were added. No statement of the manuscript
was changed. Every numbered statement, equation and section keeps its
number in the delivered PDF; that PDF is not shipped, and
\texttt{article.pdf} is a build of this text.

The delivered files are shipped under new paths:
\path{checks/check_finite.py} is \path{code/check_finite.py},
\path{checks/results.json} and \path{checks/run.txt} are
\path{data/results.json} and \path{data/run.txt} (the two are
byte-identical), and \path{build.sh} is \path{code/build.sh};
\path{PROOF_AUDIT.md}, \path{SOURCES.md} and \path{REPOSITORY_UPDATE.md}
keep their names. The delivered checksum file is not shipped.
\path{REPOSITORY_UPDATE.md} proposes new wording for the plan's C2 entry.
It is filed verbatim as a record of the delivery and has \emph{not} been
applied: the plan, its README and the synthesis are unchanged.

\subsection{\wtag\ Formal status and the Lean library}
\label{cr14:sec:wformal}
No statement of this report is formalized, in Lean or in Rocq; that it
sits in the research programme of the Lean/Rocq project
\path{Computability/TuringDegrees} gives it no formal status. The project's
Lean library \texttt{CoarseDegrees}
(\path{Computability/TuringDegrees/Lean/CoarseDegrees}) formalizes coarse,
not effective-dense, descriptions: it has no $\bot$-valued descriptions, no
$\equiv_{\ued}$ or $\equiv_{\ned}$, no eventually different or DNR
functions, no highness and no jump operator. The report relies on none of
its eight admitted statements. What it shares with the library is basic:
the predicate \path{OneGeneric} (\path{Sets.lean}, line 136) has the
definition used here, and the existence of $1$-generic sets is proved there
as \path{exists_oneGeneric} (\path{Generic.lean}, line 122), without the
bound $G\leT\emptyset'$ of Proposition~\ref{cr14:prop:low}. The external
input of the low refinement, Theorem~\ref{cr14:thm:KMS}, has no Lean
statement. Question~\ref{cr14:q:lean} and Section~\ref{cr14:sec:formal} are
a plan: no Lean or Rocq code was delivered or is shipped.

\subsection{\wtag\ Programme targets addressed, and what stays open}
\label{cr14:sec:wtargets}
Line numbers refer to
\path{Computability/TuringDegrees/Research/CoarseDegrees/research-plan/turing_degrees_unified.tex}
(``plan''), its \path{README.md}, and
\path{.../research-synthesis/Turing_Degrees_Synthesis.tex} (``synthesis''),
all identical at the pin and in the current tree.
\begin{itemize}
\item \emph{C2} (plan:563--572; plan README:20--22; synthesis:969, ``None of
the reports bears on the effective-dense question in general''; plan:1647
lists ``C2 along the line of attack'' among the first choices for fresh
work). Theorem~\ref{cr14:thm:noleast} refutes the universal C2 statement,
in both the uniform and the nonuniform version, for every $1$-generic with
$\CA(G)$: all $2$-generics directly, all non-high $1$-generics through the
external Theorem~\ref{cr14:thm:KMS}, and low $\Delta^0_2$ examples. The
plan's line of attack (plan:570) is followed in its first two steps
(Lemma~\ref{cr14:lem:equivalence} is step (i) with $0$ in place of $2$;
Corollary~\ref{cr14:cor:omissions} is step (ii)), and its step (iii), ``the
gap'', is \emph{bypassed, not proved}: the hypothesis $\CA(G)$ and
Theorem~\ref{cr14:thm:escape} replace the forcing argument the plan asks
for. For $1$-generics without $\CA$, which are exactly the high ones by
Theorem~\ref{cr14:thm:KMS} and Lemma~\ref{cr14:lem:nodnr}, step (iii) is
still open (Question~\ref{cr14:q:high}). So the C2 question is answered
negatively (unrefereed), and the gap in the plan's sketch is re-scoped to
high $1$-generics, not closed. The ``weaker, still informative result''
suggested at plan:572 is exceeded for these $G$; the positive test on the
guarded c.e.\ set suggested there is not attempted; the literature audit
suggested there was done for Astor--Hirschfeldt--Jockusch and Gerdes, and
is bounded (Appendix~\ref{cr14:app:sources}).
\item \emph{C3} (plan:574): not solved; Question~\ref{cr14:q:binary} notes a
possible connection.
\item \emph{Minimal elements and coinitiality} for effective-dense classes
are open (Questions~\ref{cr14:q:minimal} and~\ref{cr14:q:coinitial}); the
synthesis items for coarse spectra (synthesis:970, 973) and C4--C8 are not
addressed.
\item The delivered \path{REPOSITORY_UPDATE.md} proposes the corresponding
change of the plan's C2 status. It is not applied here; amending the plan
and the synthesis is separate programme work.
\end{itemize}

\section{Definitions and conventions}
\label{cr14:sec:definitions}
\subsection{Descriptions and their two reducibilities}
We identify a set with its characteristic function. All ordinary functions
under discussion are total unless explicitly called partial. A function
oracle can equivalently be coded by its graph, so ordinary Turing
reducibility between sets and functions is unambiguous.

For $S\subseteq\N$ and $N\geq1$, put
\[
 \rho_N(S)=\frac{|S\cap[0,N)|}{N},\qquad
 \Znull=\{S:\lim_{N\to\infty}\rho_N(S)=0\}.
\]
We write $S\substar T$ when $S\setminus T$ is finite. Finite unions of
members of $\Znull$ belong to $\Znull$, since prefix counts are subadditive.

\begin{definition}
An \emph{effective-dense description} of $f\in\Baire$ is a total function
$d:\N\to\N\cup\{\bot\}$ satisfying
\[
 d(n)\in\{f(n),\bot\}\quad(n\in\N),\qquad
 S_d:=\{n:d(n)=\bot\}\in\Znull.
\]
Let $\ED(f)$ be the collection of all such descriptions.
\end{definition}

The explicit output $\bot$ can be coded by a distinguished natural number,
with numerical answers shifted by one. This coding does not change degrees.
The definition agrees with the strong-domain convention for effective dense
computability.\cite[Definitions~1.4 and~3.1]{AHJ}
In particular, $f$ itself is a member of $\ED(f)$, and if $d\leT A$ then
$S_d\leT A$.

\begin{definition}
For $f,g\in\Baire$,
\[
 f\leN g\quad\Longleftrightarrow\quad
 (\forall d\in\ED(g))(\exists e\in\ED(f))\ e\leT d.
\]
We have $f\leU g$ if there is one oracle Turing functional $\Phi$ such that
$\Phi^d$ is total and belongs to $\ED(f)$ for every $d\in\ED(g)$.
Equivalence means reduction in both directions.
\end{definition}

The uniform relation implies the nonuniform relation. Reflexivity is
witnessed by copying a description. For transitivity, compose the two
functional transformations in the uniform case; in the nonuniform case,
choose a description at the intermediate step and use transitivity of
$\leT$. Thus the indicated equivalence classes are well-defined.

A key implication used repeatedly is
\begin{equation}\label{cr14:eq:representativecomputes}
 G\leN g\quad\Longrightarrow\quad
 (\exists d\in\ED(G))\ d\leT g.
\end{equation}
It follows by using $g$ itself as the input description of $g$. It does not
say that an arbitrary description of $G$, recoded as a total numerical
function, must be an equivalent representative.

\subsection{Finite oracle computations and genericity}
Let $2^{<\N}$ be the finite binary strings. The relation
$\sigma\prefix X$ means that $\sigma$ is an initial segment of $X$;
$\tau\extends\sigma$ means extension. A computation
$\Phi^\sigma(n)\downarrow=v$ is witnessed by a finite run using only
coordinates below $|\sigma|$. It persists under extensions. Finite
convergence is computably enumerable, and convergence within a specified
number of steps is decidable.

\begin{definition}
A binary set $G$ is \emph{$1$-generic} if, for every computably enumerable
$W\subseteq2^{<\N}$, either some $\sigma\prefix G$ belongs to $W$, or
some $\sigma\prefix G$ has no extension in $W$.
It is \emph{$2$-generic} if the same condition holds for every
$\Sigma^0_2$ set of strings.
\end{definition}

Avoiding a set of strings always means the stronger, cylinder-wise
alternative in this definition, not merely having no prefix in the set.
The latter alone is insufficient for the arguments below. A $2$-generic is
$1$-generic. If a c.e. set of strings is dense above a prefix of a $1$-generic,
that generic must meet it.

For fixed computable $T$, let $P_T(X)=X_T$ as in~\eqref{cr14:eq:maskintro};
for a finite string, $P_T(\sigma)$ is the same-length string obtained by
zeroing its $T$ coordinates. The map is computable, preserves extension,
and ignores the actual input bits on $T$.

\subsection{Least, minimal, and coinitial}
For a class $\mathcal C$ of oracles, $g\in\mathcal C$ has least degree in
$\mathcal C$ if $g\leT h$ for every $h\in\mathcal C$. It has minimal degree
in $\mathcal C$ if no $h\in\mathcal C$ satisfies $h\ltT g$.
Least implies minimal, but the converse need not hold. A subfamily
$\mathcal A\subseteq\mathcal C$ is \emph{coinitial} if every member of
$\mathcal C$ computes some member of $\mathcal A$.
We shall exclude finite coinitial families and certain uniformly presented
countable ones. We do not exclude arbitrary countable coinitial families.

\section{Computable erasures and the generic mask lattice}
\label{cr14:sec:erasures}
\subsection{Two transformations that never give a wrong answer}
\begin{lemma}[Computable-null erasure]\label{cr14:lem:equivalence}
For every binary $X$ and computable $T\in\Znull$, we have
$X_T\eqU X$ and $X_T\leT X$.
\end{lemma}
\begin{proof}
Given $d\in\ED(X)$, define
\[
 F_T(d)(n)=\begin{cases}0,&n\in T,\\d(n),&n\notin T.\end{cases}
\]
This is an effective-dense description of $X_T$: it is correct on $T$,
correct or $\bot$ elsewhere, and has omission set $S_d\setminus T$.
Given $e\in\ED(X_T)$, define
\[
 H_T(e)(n)=\begin{cases}\bot,&n\in T,\\e(n),&n\notin T.\end{cases}
\]
It describes $X$, with omission set $T\cup S_e$. Both sets of omissions
are null. The two transformations are single total oracle procedures with
the decidable mask hard-coded; no search for a density-convergence rate is
used. Finally, an oracle for $X$ computes $X_T$ directly.
\end{proof}

\subsection{A deleted coordinate cannot be predicted infinitely often}
\begin{lemma}[Generic nonprediction]\label{cr14:lem:prediction}
Let $G$ be $1$-generic and let $T$ be computable. Suppose a partial
$G_T$-computable function $\psi$ has the property that, whenever
$n\in T$ and $\psi(n)\downarrow$, its value is $G(n)$.
Then $T\cap\dom\psi$ is finite.
\end{lemma}
\begin{proof}
Fix an index $e$ with $\psi=\Phi_e^{G_T}$. Let $W$ be the c.e. set of
strings $\sigma$ for which there are $n\in T$ below $|\sigma|$ and
$b\in\{0,1\}$ such that
\[
 \Phi_e^{P_T(\sigma)}(n)\downarrow=b,
 \qquad \sigma(n)=1-b.
\]
No prefix of $G$ is in $W$, since such a finite computation would persist
on $G_T$ and give an incorrect prediction. By $1$-genericity, choose
$\sigma_0\prefix G$ with no extension in $W$.

Suppose $n\in T$, $n\geq|\sigma_0|$, and
$\Phi_e^{G_T}(n)\downarrow=b$. Take a sufficiently long prefix $\tau$ of
$G$ to witness this convergence and include coordinate $n$. Replace the
$n$th bit of $\tau$ by $1-b$, obtaining $\widehat\tau$. This still extends
$\sigma_0$, and
$P_T(\widehat\tau)=P_T(\tau)$ because $n\in T$. Thus
$\widehat\tau\in W$, a contradiction. All convergent inputs in $T$ are
therefore below $|\sigma_0|$.
\end{proof}

\begin{corollary}[Forced omissions]\label{cr14:cor:omissions}
If $G$ is $1$-generic, $T$ is computable, and
$d\in\ED(G)$ satisfies $d\leT G_T$, then $T\substar S_d$.
\end{corollary}
\begin{proof}
On input $n$, compute $d(n)$ and return its value if numerical; otherwise
run forever. This gives a partial $G_T$-computable predictor that is always
correct. Its convergent inputs in $T$ are exactly $T\setminus S_d$.
Apply Lemma~\ref{cr14:lem:prediction}.
\end{proof}

\subsection{An exact order calculation}
\begin{theorem}[Mask order]\label{cr14:thm:maskorder}
For a $1$-generic $G$ and computable $C,D\subseteq\N$,
\begin{equation}\label{cr14:eq:maskorder}
 G_C\leT G_D\quad\Longleftrightarrow\quad D\substar C.
\end{equation}
Consequently $G_C\eqT G_D$ exactly when $C\symdiff D$ is finite.
\end{theorem}
\begin{proof}
If $D\setminus C$ is finite, compute $G_C$ from $G_D$ by outputting zero
on $C$, copying $G_D$ off $C\cup D$, and hard-coding the finitely many
correct bits on $D\setminus C$.

Conversely, suppose $G_C\leT G_D$ and $D\setminus C$ is infinite.
On that infinite computable set, $G_C$ supplies the original bits of $G$.
It would therefore be a partial $G_D$-computable predictor correct on an
infinite subset of $D$, contrary to Lemma~\ref{cr14:lem:prediction}.
\end{proof}

In particular, every infinite computable mask strictly lowers the degree
of a $1$-generic. If $C\subseteq D$ and $D\setminus C$ is infinite, then
$G_D\ltT G_C$. This does not require $G_C$ itself to be generic.

\begin{theorem}[Finite joins and actual greatest lower bounds]
\label{cr14:thm:masklattice}
For computable $C,D$ and a $1$-generic $G$,
\[
 G_C\oplus G_D\eqT G_{C\cap D},
 \qquad
 \Low(G_C)\cap\Low(G_D)=\Low(G_{C\cup D}),
\]
where $\Low(X)=\{A:A\leT X\}$. Thus the second identity gives a greatest
lower bound in the full Turing degrees, not merely inside the mask family.
\end{theorem}
\begin{proof}
The join identity is direct. Off $C\cap D$, at least one oracle retains
the original bit; decidability of $C,D$ tells us which one to use.
Conversely, $G_{C\cap D}$ computes both masked oracles by further zeroing.
Also $G_{C\cup D}$ computes from either, so one inclusion in the lower-cone
identity is immediate.

For the other, fix a total function
$A=\Phi^{G_C}=\Psi^{G_D}$. Let $W$ be the c.e. set of strings $\sigma$
witnessing, at some input $n$,
\[
 \Phi^{P_C(\sigma)}(n)\downarrow\ne
       \Psi^{P_D(\sigma)}(n)\downarrow.
\]
No prefix of $G$ lies in $W$. Choose $\sigma_0\prefix G$ such that no
extension is in $W$. Put $H=G_{C\cup D}$.

To compute $A(n)$ from $H$, search for a finite string $\tau$ extending
$\sigma_0$ which agrees with $H$ at all coordinates outside $C\cup D$
below $|\tau|$, and for which
$\Phi^{P_C(\tau)}(n)\downarrow=v$. The compatibility test is decidable
in $H$; the search halts because a sufficiently long actual prefix of $G$
is a witness.

We claim that every value found equals $A(n)$. Take a long actual prefix
$\eta\prefix G$, extending $\sigma_0$, which witnesses
$\Psi^{P_D(\eta)}(n)\downarrow=A(n)$. The private coordinates used by the
first oracle are in $D\setminus C$; those used by the second are in
$C\setminus D$. They are disjoint. On the shared coordinates outside
$C\cup D$, the two strings agree by the compatibility test. On
$C\cap D$, both projected oracles are zero. We may therefore amalgamate
the two finite computations into one string $\nu\extends\sigma_0$,
preserving the first computation from $\tau$ and the second from $\eta$.
If $v\ne A(n)$, this gives $\nu\in W$, a contradiction. The search
computes $A$ from $H$.
\end{proof}

It follows that
\[
 (\Comp\cap\Znull)/\Fin\ \longrightarrow\ \Spec([G]_{\ued}),
 \qquad [C]\longmapsto\degT(G_C),
\]
is an order-reversing injection with union corresponding to meet and
intersection to join. This is a precise lattice structure \emph{inside one}
uniform effective-dense class. It has no least member: any computable null
$C$ can be enlarged by an infinite computable null subset of its complement.
For example, find successively a point outside $C$ beyond $2^n$ and the
previous chosen point. Theorem~\ref{cr14:thm:maskorder} then supplies a strictly
smaller member. This last observation by itself does not rule out a least
representative lying outside the mask family; that is the purpose of the
next sections.

\section{Escaping a null set with a computable sparse selector}
\label{cr14:sec:escape}
\subsection{Packing prescribed capacities into every-prefix budgets}
\begin{lemma}[Budgeted blocks]\label{cr14:lem:blocks}
Let $b:\N\to\N$ be computable, nondecreasing, and unbounded, and let
$(a_n)$ be a computable sequence of positive integers. Put
$Q_n=\sum_{j\leq n}a_j$. There is a computable sequence $(m_n)$ such that
\begin{gather}
 m_n\geq4^{n+1},\qquad m_{n+1}\geq4m_n,\qquad m_n\geq2n+3,
 \label{cr14:eq:blockspacing}\\
 m_n\geq(n+1)Q_n,\qquad b(m_n)\geq Q_n.\label{cr14:eq:blockbudget}
\end{gather}
Let $I_n=[m_n,2m_n)$. Any set $T$ contained in $\bigcup_n I_n$ with
$|T\cap I_n|\leq a_n$ satisfies
\[
 |T\cap[0,N)|\leq b(N)\quad(N\in\N),\qquad T\in\Znull.
\]
No computability of $T$ is needed for this last conclusion.
\end{lemma}
\begin{proof}
Find $m_n$ successively by an ordinary search satisfying the finitely many
lower bounds and $b(m_n)\geq Q_n$. Unboundedness of $b$ makes each search
halt. The intervals are disjoint by~\eqref{cr14:eq:blockspacing}.
For $N<m_0$ the count is zero. If $m_n\leq N<m_{n+1}$, then
\[
 |T\cap[0,N)|\leq Q_n\leq b(m_n)\leq b(N).
\]
Moreover, for $N\geq1$ in the same range,
\[
 \rho_N(T)\leq Q_n/m_n\leq1/(n+1).
\]
As $N$ tends to infinity, so does $n$, proving density zero. Notice that
this estimate covers prefixes \emph{inside} a block as well as endpoints.
\end{proof}

\subsection{One missed point per interval}
\begin{definition}
For any sequence of blocks as in Lemma~\ref{cr14:lem:blocks} and any set $S$,
define
\begin{equation}\label{cr14:eq:choice}
 f_S(n)=
 \begin{cases}
 \min\{k<m_n:m_n+k\notin S\},&I_n\nsubseteq S,\\
 0,&I_n\subseteq S.
 \end{cases}
\end{equation}
For a total computable $h:\N\to\N$, define its selector
\begin{equation}\label{cr14:eq:selector}
 T_h=\{m_n+(h(n)\bmod m_n):n\in\N\}.
\end{equation}
\end{definition}

The function $f_S$ is total and computable from $S$. It requires only a
finite search on each input, even if $m_n$ is very large. The set $T_h$ is
computable: find the unique block which could contain the proposed input,
if one exists, and compute the one selected point in that block. It is not
merely c.e. The total running time of $h(n)$ need not have been predicted
when the interval endpoints were chosen.

\begin{lemma}[Null covers imply eventual difference]\label{cr14:lem:cover}
Suppose $S\in\Znull$. For all sufficiently large $n$, $I_n\nsubseteq S$.
If $T_h\substar S$ for every total computable $h$, then $f_S$ is eventually
different from every total computable function.
\end{lemma}
\begin{proof}
If infinitely many $I_n$ were contained in $S$, then at the corresponding
endpoints $\rho_{2m_n}(S)\geq1/2$, contradicting density zero.
Fix computable $h$. For all sufficiently large $n$, the chosen point of
$T_h$ belongs to $S$, whereas $m_n+f_S(n)$ is outside $S$.
If $f_S(n)=h(n)$, then $h(n)<m_n$ and the two points coincide, a
contradiction. Thus $f_S(n)\ne h(n)$ eventually.
\end{proof}

\begin{theorem}[Budgeted sparse escape]\label{cr14:thm:escape}
If $\CA(G)$, $S\leT G$, and $S\in\Znull$, then for every unbounded
computable order $b$ there is an infinite computable $T\in\Znull$ with
\begin{equation}\label{cr14:eq:escapebudget}
 |T\cap[0,N)|\leq b(N)\quad(N\in\N),\qquad |T\setminus S|=\infty.
\end{equation}
\end{theorem}
\begin{proof}
Use Lemma~\ref{cr14:lem:blocks} with $a_n=1$. The total function $f_S$ of
\eqref{cr14:eq:choice} is computable from $G$. By $\CA(G)$, some total
computable $h$ agrees with it infinitely often. On all sufficiently large
agreement inputs, the corresponding point in $T_h$ is a genuine point
outside $S$. Thus $T=T_h$ has infinitely many points outside $S$.
It has exactly one point in each block, so it is infinite, is computable,
and satisfies the budget and density conclusions of Lemma~\ref{cr14:lem:blocks}.
\end{proof}

\begin{remark}[What happened to the timing gap?]\label{cr14:rem:timing}
The argument never attempts to enumerate a sparse set of potential witnesses
and then infer that the enumeration has a computable membership test. It
first uses $\CA(G)$ to obtain a \emph{total computable function} $h$, and
then defines exactly one point in each computably located interval.
Unbounded convergence times cause no difficulty for a membership algorithm
which is permitted to wait for a guaranteed total computation. This is the
specific replacement for the delayed-witness construction discussed in C2.
\end{remark}

A useful independent consequence is that any oracle computing a null set
which almost contains every computable null set must compute an eventually
different function. Indeed all the selectors in~\eqref{cr14:eq:selector} are
computable and null. The choice function is even bounded by the computable
function $m_n$. This observation explains why the existence of abstract
null upper bounds does not settle an effective question about them.

\section{Classical sources of computable agreement}
\label{cr14:sec:agreement}
\subsection{A direct proof for every \texorpdfstring{$2$}{2}-generic}
The following fact is classical; it is consistent with the formulation in
terms of weak meager engulfing recorded by Brendle--Brooke-Taylor--Ng--Nies.
We supply the finite-extension proof needed here.\cite[Theorem~4.2 and
Section~4, item~(5)]{BBNN}

\begin{theorem}\label{cr14:thm:twoGenericCA}
Every $2$-generic $G$ satisfies $\CA(G)$.
\end{theorem}
\begin{proof}
Fix a functional $\Gamma$ with $f=\Gamma^G$ total. Consider
\[
 B_\Gamma=\{\sigma:(\exists n)(\forall\tau\extends\sigma)
                         \Gamma^\tau(n)\uparrow\}.
\]
Here divergence of a finite-oracle computation means that no finite-time
convergence witness exists. The displayed condition has the form
$\exists n\,\forall(\tau,s)$ with a decidable matrix, so $B_\Gamma$ is
$\Sigma^0_2$. No prefix of $G$ is in it, because $\Gamma^G$ is total.
There is therefore $\sigma_0\prefix G$ such that no extension of
$\sigma_0$ is in $B_\Gamma$. Equivalently,
\begin{equation}\label{cr14:eq:densetotality}
 (\forall\tau\extends\sigma_0)(\forall n)
       (\exists\nu\extends\tau)\;\Gamma^\nu(n)\downarrow.
\end{equation}

Choose a computable sequence $(\tau_n)$ listing every finite extension of
$\sigma_0$ infinitely often. On input $n$, search for an extension
$\nu\extends\tau_n$ with $\Gamma^\nu(n)\downarrow$, and set $h(n)$ to
the value found. Equation~\eqref{cr14:eq:densetotality} makes $h$ total and
computable. The finite string $\sigma_0$ is hard-coded; it need not be
computably selected from an index for $G$.

For each $k$, the set
\[
 U_k=\{\nu:(\exists n\geq k)\;\Gamma^\nu(n)\downarrow=h(n)\}
\]
is c.e. and dense above $\sigma_0$. To see density, take any extension
$\tau$ of $\sigma_0$, choose $n\geq k$ with $\tau_n=\tau$, and use the
extension found while defining $h(n)$. Since $G$ is also $1$-generic, it
meets $U_k$. Thus $f(n)=h(n)$ for some $n\geq k$, for every $k$.
\end{proof}

This theorem, together with the preceding sections and
Section~\ref{cr14:sec:mainproof}, gives a proof of the negative least-degree
answer without importing a highness classification or a basis theorem.

\subsection{Non-high \texorpdfstring{$1$}{1}-generics and low examples}
A total function $d:\N\to\N$ is \emph{diagonally nonrecursive}, or DNR,
if $d(e)\ne\varphi_e(e)$ whenever the latter is defined. An oracle has
DNR degree if it computes such a function. It is \emph{high} if its jump
computes $\doublezero$.

\begin{theorem}[Classical characterization; external input]
\label{cr14:thm:KMS}
An oracle computes a function eventually different from every total
computable function if and only if it is high or has DNR degree.
\end{theorem}

This is Kjos-Hanssen--Merkle--Stephan, Theorem~5.1.\cite[Theorem~5.1]{KMS}
We use its stated equivalence, not any of its Kolmogorov-complexity
characterizations. The theorem is not a result of the present paper.

\begin{lemma}[A standard genericity consequence]\label{cr14:lem:nodnr}
No $1$-generic computes a DNR function.
\end{lemma}
\begin{proof}
Suppose $d=\Gamma^G$ is total and DNR. The set
\[
 W=\{\sigma:(\exists e)\ \Gamma^\sigma(e)\downarrow
                                  =\varphi_e(e)\downarrow\}
\]
is c.e. and has no prefix of $G$. Choose a prefix $\sigma_0$ of $G$ with
no extension in $W$. Define $h(e)$ by searching for any extension
$\tau\extends\sigma_0$ with $\Gamma^\tau(e)\downarrow$ and returning
the value. Since $\Gamma^G$ is total, this is a total computable function.
The choice of $\sigma_0$ guarantees that $h$ is DNR. But if $j$ is a
computable index for $h$, then $\varphi_j(j)=h(j)$, a contradiction.
\end{proof}

\begin{corollary}\label{cr14:cor:nonhigh}
Every non-high $1$-generic has $\CA$. In particular, every low $1$-generic
has $\CA$.
\end{corollary}
\begin{proof}
By Lemma~\ref{cr14:lem:nodnr} neither alternative in
Theorem~\ref{cr14:thm:KMS} holds for a non-high $1$-generic.
A low oracle is not high, by strictness of the Turing jump.
\end{proof}

For completeness, the low examples needed in the main theorem can be
obtained explicitly relative to the halting set.

\begin{proposition}[Low $\Delta^0_2$ generic witnesses]\label{cr14:prop:low}
There is a $1$-generic $G\leT\jumpzero$. Every $1$-generic below
$\jumpzero$ is low and noncomputable.
\end{proposition}
\begin{proof}
Enumerate all c.e. sets of finite strings as $(W_e)$. Starting with the
empty string, at stage $e$ use $\jumpzero$ to decide whether the current
string has an extension in $W_e$. If yes, find such an extension by
ordinary enumeration and adopt it; if not, keep the current string.
Append a zero to ensure growth. The union $G$ is computable from
$\jumpzero$ and meets or avoids each $W_e$.

For any $1$-generic $G$, membership in $G'$ can be decided from
$G\oplus\jumpzero$. For the computation $\Phi_e^G(e)$, search along
prefixes of $G$ for either a convergence witness or a prefix having no
extension on which the computation converges. The latter test is decidable
in $\jumpzero$. Genericity guarantees one of these alternatives.
Hence $G'\leT G\oplus\jumpzero$. If $G\leT\jumpzero$, it follows that
$G'\eqT\jumpzero$. Finally, a computable $G$ would be a predictor of its
own bits from the all-zero oracle $G_\N$, violating
Lemma~\ref{cr14:lem:prediction} with $T=\N$.
\end{proof}

The generalized-low calculation in this proof is standard; it is also
recalled in the discussion surrounding Proposition~5.5 of
Astor--Hirschfeldt--Jockusch.\cite{AHJ}
For our application the important point is the separation of dependencies:
the direct $2$-generic route is elementary, whereas the sharper low-witness
route uses Theorem~\ref{cr14:thm:KMS} in addition.

\section{No least representative, even under an arbitrary sparse budget}
\label{cr14:sec:mainproof}
\subsection{The local theorem}
\begin{theorem}\label{cr14:thm:local}
Let $G$ be $1$-generic with $\CA(G)$. Let $A\leT G$ and suppose some
$d\in\ED(G)$ is computable from $A$. For any unbounded computable order
$b$, there is an infinite computable mask $T$ satisfying all conclusions of
Theorem~\ref{cr14:thm:localintro}.
\end{theorem}
\begin{proof}
Let $S=S_d$. Then $S\in\Znull$ and $S\leT A\leT G$. By
Theorem~\ref{cr14:thm:escape}, choose an infinite computable null set $T$
obeying the budget and with $T\setminus S$ infinite.
If $A\leT G_T$, then $d\leT G_T$, so
Corollary~\ref{cr14:cor:omissions} would imply $T\substar S$.
This contradiction proves $A\nleT G_T$.

Lemma~\ref{cr14:lem:equivalence} gives $G_T\eqU G$ and $G_T\leT G$.
The inequality is strict by Theorem~\ref{cr14:thm:maskorder}, since $T$ is
infinite. Finally $G\symdiff G_T\subseteq T$, so the same budget bounds
the actual disagreements.
\end{proof}

\subsection{Checking the quantifiers of the original question}
\begin{theorem}[Negative answer to C2]\label{cr14:thm:noleast}
If $G$ is $1$-generic and has $\CA(G)$, neither $[G]_{\ned}$ nor
$[G]_{\ued}$ has an element of least Turing degree, even if only binary
representatives are allowed.
\end{theorem}
\begin{proof}
Take any proposed representative $g\in[G]_{\ned}$.
If $g\nleT G$, the representative $G$ already witnesses that $g$ is not
least. In fact every $G_T\leT G$ also fails to compute $g$.

Otherwise $g\leT G$. Since $G\leN g$, implication
\eqref{cr14:eq:representativecomputes} supplies $d\in\ED(G)$ with $d\leT g$.
Apply Theorem~\ref{cr14:thm:local} with $A=g$. The resulting binary
representative $G_T\eqU G$ does not compute $g$. It belongs to the
nonuniform class as well, so $g$ is not least there.

If $g$ was proposed as a least member only of $[G]_{\ued}$, the same
argument applies: this class is contained in $[G]_{\ned}$, and both
comparison witnesses, $G$ and $G_T$, are in the uniform class. They are
binary, so the proof also excludes a least degree in either binary
subclass.
\end{proof}

Together with Corollary~\ref{cr14:cor:nonhigh} and
Proposition~\ref{cr14:prop:low}, this proves Theorem~\ref{cr14:thm:mainintro},
including the low $\Delta^0_2$ refinement.

\begin{corollary}[A countable test family already has no in-class lower bound]
\label{cr14:cor:testfamily}
Fix $G$ as above and an unbounded computable order $b$. Let
\[
 \mathcal M_b(G)=\{G_T:T\text{ computable},\ T\in\Znull,
                      \ (\forall N)\ |T\cap[0,N)|\leq b(N)\}.
\]
No member of $[G]_{\ned}$ is a Turing lower bound of
$\mathcal M_b(G)$. In particular the nonuniform class has no least member
even when leastness is tested only against this countable family of binary
representatives in the uniform class.
\end{corollary}
\begin{proof}
A lower bound $g$ is below $G$, since the empty mask is allowed. As above,
$g$ computes a description of $G$. Theorem~\ref{cr14:thm:local} supplies a
member of $\mathcal M_b(G)$ not computing $g$.
\end{proof}

\subsection{Finite simultaneous avoidance and a sharp budget boundary}
\begin{corollary}[Finite targets]\label{cr14:cor:finite}
For every finite family $g_0,\ldots,g_{k-1}\in[G]_{\ned}$ and every
unbounded computable order $b$, one computable null mask $T$ within budget
satisfies $g_i\nleT G_T$ for all $i<k$. Consequently neither the uniform
nor the nonuniform class has a finite coinitial subset.
\end{corollary}
\begin{proof}
Any $g_i\nleT G$ is automatically not below $G_T$. For each remaining
$i$, choose $d_i\in\ED(G)$ with $d_i\leT g_i\leT G$.
The finite union $S=\bigcup_i S_{d_i}$ is null and computable from $G$.
Use Theorem~\ref{cr14:thm:escape} to make $T\setminus S$ infinite. Then
$T\setminus S_{d_i}$ is infinite for each relevant $i$, and
Corollary~\ref{cr14:cor:omissions} rules out $g_i\leT G_T$.
If there are no relevant targets, any infinite budgeted selector suffices.
An equivalent representative not computing any member of the proposed
finite family disproves coinitiality.
\end{proof}

\begin{proposition}[Why unboundedness is necessary]\label{cr14:prop:bounded}
If $|T\cap[0,N)|\leq M$ for all $N$, then $T$ is finite and
$G_T\eqT G$. Such masks cannot witness $A\nleT G_T$ for any $A\leT G$.
\end{proposition}
\begin{proof}
An infinite $T$ would have a prefix containing more than $M$ elements.
Finite changes preserve ordinary Turing degree by hard-coding the changed
bits.
\end{proof}

Thus the distinction between bounded and unbounded budgets is exact for
this local erasure method. A fixed-budget neighborhood is not being called
an equivalence class: budgets need not be preserved under composition.

\section{Uniformly presented countable families}
\label{cr14:sec:countable}
\subsection{An effective version of the null-set P-ideal calculation}
A countable family of null sets has a null upper bound modulo finite sets.
For the present application the computational presentation of the family
matters. The following direct construction needs a uniform membership
procedure, but needs no density-convergence moduli and no extra jump.

\begin{lemma}[Uniform null upper bounds]\label{cr14:lem:Pideal}
Suppose $(S_i)_{i\in\N}$ is uniformly computable from an oracle $G$ and
each $S_i$ is in $\Znull$. There is $S\leT G$ in $\Znull$ with
$S_i\substar S$ for every $i$.
\end{lemma}
\begin{proof}
Use the dyadic blocks $J_n=[2^n,2^{n+1})$. For $k\geq1$, define
\[
 r_k(n)=\frac{|(\bigcup_{i<k}S_i)\cap J_n|}{2^n}.
\]
For each fixed $k$, $r_k(n)\to0$, since the numerator is at most
$|\bigcup_{i<k}S_i\cap[0,2^{n+1})|$. Let
\[
 k(n)=\max\bigl(\{0\}\cup
     \{k:1\leq k\leq n+1\text{ and }r_k(n)\leq1/k\}\bigr).
\]
All tests here are finite rational comparisons, uniformly decidable in
$G$. For each fixed $k$, that $k$ is eligible at all sufficiently large
$n$. Hence $k(n)\to\infty$. Set
\[
 S\cap J_n=\left(\bigcup_{i<k(n)}S_i\right)\cap J_n,
 \qquad 0\notin S.
\]
This defines $S\leT G$. Since $k(n)>i$ eventually, $S_i\substar S$.
The fraction of $S$ in block $J_n$ is at most $1/k(n)$ when $k(n)>0$,
and is zero otherwise; it therefore tends to zero.

To pass from block fractions to full density, fix $\varepsilon>0$ and
choose $n_0$ so that all later block fractions are at most $\varepsilon$.
For $2^n\leq N<2^{n+1}$ with $n\geq n_0$,
\[
 |S\cap[0,N)|\leq 2^{n_0}
       +\varepsilon\sum_{j=n_0}^{n}2^j
 \leq 2^{n_0}+\varepsilon 2^{n+1}.
\]
Thus $\limsup_N\rho_N(S)\leq2\varepsilon$.
As $\varepsilon$ was arbitrary, $S\in\Znull$.
\end{proof}

\subsection{One computable mask defeating countably many oracles}
\begin{theorem}[Uniform countable simultaneous avoidance]
\label{cr14:thm:countable}
Let $G$ be $1$-generic with $\CA(G)$. Suppose $(d_i)_{i\in\N}$ is a
uniformly $G$-computable sequence, each $d_i\in\ED(G)$. Let $A_i$ be
any oracles with $d_i\leT A_i$ for every $i$; the latter reductions need
not be uniform. For every unbounded computable order $b$, there is an
infinite computable null mask $T$ within budget such that
\[
 (\forall i)\quad A_i\nleT G_T.
\]
\end{theorem}
\begin{proof}
The omission sets $S_i=S_{d_i}$ are uniformly $G$-computable and null.
Lemma~\ref{cr14:lem:Pideal} supplies a null $S\leT G$ with $S_i\substar S$
for every $i$. By Theorem~\ref{cr14:thm:escape}, choose a computable budgeted
$T$ with $T\setminus S$ infinite. For each $i$, only finitely many of
these escaping points can belong to $S_i\setminus S$, so
$T\setminus S_i$ is infinite. If $A_i\leT G_T$, then
$d_i\leT G_T$, contradicting Corollary~\ref{cr14:cor:omissions}.
\end{proof}

An immediate special case is that no uniformly $G$-computable family of
effective-dense descriptions has the property that every computable-null
erasure of $G$ computes one of those descriptions. Theorem~\ref{cr14:thm:countable}
also excludes a countable coinitial family of representatives whenever
suitable descriptions below its members have such a uniform presentation.

\begin{remark}[The quantifier that cannot be discarded]\label{cr14:rem:uniform}
A sequence of sets individually computable from $G$ need not be uniformly
computable from $G$. Likewise, the assertion that each proposed
representative computes some description does not uniformly select indices
for correct total descriptions. Lemma~\ref{cr14:lem:Pideal} makes finitely many
queries across the first $k$ rows on each block; it needs the uniform
procedure for those rows. This paper therefore does not infer the absence
of arbitrary countable coinitial families from Theorem~\ref{cr14:thm:countable}.
\end{remark}

\section{A packed descending chain with no in-class lower bound}
\label{cr14:sec:chain}
\subsection{A nonuniform enumeration is used openly}
Fix an unbounded computable order $b$. Choose blocks $I_n=[m_n,2m_n)$
from Lemma~\ref{cr14:lem:blocks} with
\[
 a_n=2(n+1),\qquad Q_n=(n+1)(n+2).
\]
Let $(h_i)_{i\in\N}$ list all total computable functions. This is a
set-theoretic listing, \emph{not} a computable enumeration of total
computable functions. For instance, list increasingly the indices in
$\mathrm{TOT}$ and take their computed functions. Membership in
$\mathrm{TOT}$ is decidable in $\doublezero$, so such a listing can be
chosen $\doublezero$-computable at the level of indices. For each fixed
$i$, the finite list $h_0,\ldots,h_i$ consists of ordinary computable
functions whose indices may be hard-coded.

Define for all $i,n$
\[
 p_i(n)=h_i(n)\bmod m_n.
\]
For $n\geq i$, define recursively in $i$
\begin{equation}\label{cr14:eq:fresh}
 q_i(n)=\min\left([0,m_n)\setminus
      \left(\{p_i(n)\}\cup\{p_j(n),q_j(n):j<i\}\right)\right).
\end{equation}
At most $2i+1$ offsets are forbidden, while $m_n\geq2n+3>2i+1$.
Thus this value exists. For fixed $i$, it is computable on $n\geq i$.
Values on smaller $n$ are unused and may be set to zero.

Let $C_i$ be empty outside the blocks and set
\begin{equation}\label{cr14:eq:chainmask}
 C_i\cap I_n=
 \{m_n+p_j(n),\ m_n+q_j(n):0\leq j\leq\min(i,n)\}.
\end{equation}
Each $C_i$ is computable. To decide membership, locate the relevant block
and run only the finitely many computable functions with indices at most
$i$. There is no assertion that $i\mapsto$ an ordinary computable index
for $C_i$ is itself computable.

\subsection{The geometry and strict descent}
\begin{lemma}\label{cr14:lem:chaingeometry}
The masks satisfy $C_i\subseteq C_{i+1}$ and
$|C_{i+1}\setminus C_i|=\infty$. Their union
$C_\infty=\bigcup_i C_i$ is null and obeys the budget $b$ at every prefix.
\end{lemma}
\begin{proof}
Inclusion is immediate from~\eqref{cr14:eq:chainmask}. For each $n\geq i+1$,
the point $m_n+q_{i+1}(n)$ belongs to $C_{i+1}$ but not to $C_i$, by
its freshness requirement. These points are distinct in different blocks.
There are at most $2(n+1)$ points of $C_\infty$ in block $I_n$, so
Lemma~\ref{cr14:lem:blocks} gives both its nullity and the budget. Every $C_i$
inherits them. The union need not be computable.
\end{proof}

\begin{theorem}[Descending chain theorem]\label{cr14:thm:chain}
Let $G$ be $1$-generic with $\CA(G)$. For the masks above, $G_i=G_{C_i}$
satisfies all conclusions of Theorem~\ref{cr14:thm:chainintro}.
\end{theorem}
\begin{proof}
Lemma~\ref{cr14:lem:equivalence} gives $G_i\eqU G$ and $G_i\leT G$.
Theorem~\ref{cr14:thm:maskorder} and
Lemma~\ref{cr14:lem:chaingeometry} give $G_{i+1}\ltT G_i$. The budget bounds
actual disagreements because $G\symdiff G_i\subseteq C_i$.

Suppose, toward a contradiction, that $g\in[G]_{\ned}$ and
$g\leT G_i$ for every $i$. Then $g\leT G$. Choose
$d\in\ED(G)$ with $d\leT g$, and let $S=S_d$. We have $S\leT G$ and
$S\in\Znull$. For every $i$, $d\leT G_{C_i}$, so
Corollary~\ref{cr14:cor:omissions} gives
\begin{equation}\label{cr14:eq:covereachmask}
 C_i\substar S.
\end{equation}
For the particular total computable function $h_i$, the selector
$T_{h_i}$ from~\eqref{cr14:eq:selector} has all of its points on blocks $n\geq i$
in $C_i$. Thus~\eqref{cr14:eq:covereachmask} implies $T_{h_i}\substar S$.
Every total computable function appears in the list, so
Lemma~\ref{cr14:lem:cover} makes $f_S\leT G$ eventually different from every
computable function. This contradicts $\CA(G)$.

If $G$ is low, monotonicity of the jump yields
$\jumpzero\leT G_i'\leT G'\eqT\jumpzero$.
Hence each $G_i$ is low.
\end{proof}

All these degrees are noncomputable. For example,
Theorem~\ref{cr14:thm:maskorder}, with $D=\N$, would make
$G_{C_i}\leT G_\N$ imply $\N\substar C_i$, impossible for a null mask.
Thus, choosing the low generic from Proposition~\ref{cr14:prop:low}, the
construction gives an infinite strictly descending chain of nonzero low
Turing degrees in one uniform effective-dense class, without an in-class
lower bound even in the larger nonuniform class.

\subsection{Why the countable union does not give a counterargument}
The union $C_\infty$ in Lemma~\ref{cr14:lem:chaingeometry} really is null; the
argument does not exploit a failure of closure under arbitrary unions.
But $C_\infty\nleT G$. Otherwise its containment of each selector tail
would, by Lemma~\ref{cr14:lem:cover}, again give an eventually different
function below $G$. Consequently one cannot apply the computable-mask
transformation to $C_\infty$.

There is no contradiction between this construction and
Lemma~\ref{cr14:lem:Pideal}. The individual masks are computable, but the
chosen listing of them is not asserted to be uniformly $G$-computable.
The extra fresh points in~\eqref{cr14:eq:fresh} are also important: selector
functions can coincide infinitely often, so merely adjoining the next
selector need not make each consecutive degree drop strict.

For a low $G\leT\jumpzero$, the entire displayed family of masks and
masked sets has a uniform $\doublezero$ presentation using the index
listing through $\mathrm{TOT}$. This is an upper bound on the presentation,
not a claim that the bound is optimal.

\section{Finite degree patterns and interpretation of the obstruction}
\label{cr14:sec:interpretation}
\subsection{Boolean intervals inside an arbitrarily small neighborhood}
\begin{corollary}[Finite Boolean patterns]\label{cr14:cor:boolean}
Let $G$ be $1$-generic, let $b$ be an unbounded computable order, and let
$k\geq1$. Inside $[G]_{\ued}$ there are $2^k$ distinct ordinary degrees
forming a Boolean lattice under the actual Turing joins and meets.
All representatives can be chosen below $G$ and within disagreement budget
$b$ of $G$. If $G$ is low, all these degrees are low.
\end{corollary}
\begin{proof}
Choose an infinite computable null set $T$ within budget using one point per
block. Split its computable increasing enumeration into $k$ infinite
computable pieces $T_0,\ldots,T_{k-1}$ by the enumeration index modulo $k$.
For $u\subseteq\{0,\ldots,k-1\}$, let $C_u=\bigcup_{i\in u}T_i$.
All masks are subsets of $T$. Distinct $u$ give masks different on an
infinite set, so their degrees are distinct by Theorem~\ref{cr14:thm:maskorder}.
That theorem reverses inclusion, and Theorem~\ref{cr14:thm:masklattice} identifies
union with meet and intersection with join. This is the dual Boolean
lattice, which is itself Boolean. The low assertion follows by jump
monotonicity as before.
\end{proof}

For $k=2$, the two intermediate degrees are incomparable, their join is
$\degT(G)$, and their meet is $\degT(G_T)$. The meet is nonzero because
$T$ is null. Thus these are not ordinary minimal pairs. The finite geometry
is compatible with the complete lack of a least degree in the full class
under the stronger hypotheses of Theorem~\ref{cr14:thm:noleast}.

\subsection{The repository gap, resolved at the required quantifier level}\label{cr14:sec:gap}
The plan's C2 argument would make the omission set $S$ of a least
representative almost contain every computable null set. The proposed
contradiction there was not a purely set-theoretic one: null sets form a
P-ideal, so abstract upper bounds exist.\cite{ProveIt}
The present proof locates the missing computational obstruction:
\[
 \begin{gathered}
 S\leT G,\quad S\in\Znull,\quad
 (\forall T\text{ computable null})\ T\substar S\\[2pt]
 \Longrightarrow\quad
 \text{$G$ computes an eventually different function.}
 \end{gathered}
\]
For $2$-generics, or non-high $1$-generics, that consequence is impossible.
The general least-representative assertion is universal in the starting
function, so this subclass is already sufficient for a negative answer.
We have not silently completed the original sketch for every $1$-generic.

\noindent\wtag\ In the plan's terms (\path{turing_degrees_unified.tex},
line 570): step (i) is Lemma~\ref{cr14:lem:equivalence}, with the value $0$
instead of $2$ on the mask; step (ii) is
Corollary~\ref{cr14:cor:omissions}; step (iii) is replaced by the
hypothesis $\CA(G)$, which fails exactly for the high $1$-generics. The
title of this subsection should therefore be read as ``bypassed for
$1$-generics with $\CA$'': the plan's gap is re-scoped to the high
$1$-generics of Question~\ref{cr14:q:high}, not closed. The C2 question
itself, being universal, is answered negatively.

Two logically distinct facts do the work. Genericity forbids a retained
oracle from making infinitely many error-free predictions on deleted
coordinates. Computable agreement prevents a $G$-computable null omission
set from swallowing all computable sparse selectors. Neither fact alone is
the least-representative theorem.

\subsection{Six boundaries of the conclusion}
The following distinctions prevent overstatements of what has been proved.

\paragraph{Descriptions versus representatives.}
A description is allowed to contain $\bot$; a representative is a numerical
function in the equivalence class. The proof extracts a description from a
representative by~\eqref{cr14:eq:representativecomputes}. It never identifies
the entire representative spectrum with an upward closure of description
degrees.

\paragraph{Explicit omission versus an erroneous answer.}
The erased sets use numerical zero, but the reverse description reduction
\emph{explicitly} replaces every masked coordinate by $\bot$. No wrong
answer is reclassified as an omission after the fact.

\paragraph{Least versus minimal.}
The mask family has no minimal element, and the full class has no least
element. The existence of a minimal element of the \emph{full} class is
not settled by these results. A representative outside the displayed
family need not have a smaller displayed representative below it.

\paragraph{Finite versus countable uniformity.}
Finite unions of omission sets are $G$-computable when the constituent
sets are. Countable unions do not enjoy that inference without a uniform
presentation and a null-upper-bound construction. Section~\ref{cr14:sec:countable}
states exactly the uniform hypothesis it uses.

\paragraph{Actual degree order versus asymptotic degree order.}
The strict inequalities in the chain are ordinary Turing inequalities.
All chain members are equivalent for uniform effective-dense reducibility.
There is no strict chain of effective-dense degrees in that assertion.

\paragraph{Provability versus priority and formal certification.}
The arguments constitute a conventional proof submission. They are not
results certified by finite experiments, by repository inclusion, by peer
review, or by a proof assistant. The low sharpening has the named classical
input of Theorem~\ref{cr14:thm:KMS}; it is not an unproved assumption introduced
to obtain the conclusion.

\section{Further research questions and concrete next targets}
\label{cr14:sec:questions}
Except where linked to a cited question, the following are questions raised
by this development, not claims that a comprehensive search has established
their open status. Their formulations separate plausible next lemmas from
large classification problems.

\subsection{Genericity and computational strength}
\begin{question}[The high $1$-generic case]\label{cr14:q:high}
Does $[G]_{\ned}$ have no least Turing degree for every high $1$-generic $G$?
Does the same hold for $[G]_{\ued}$?
\end{question}
The present argument handles all non-high $1$-generics. In the remaining
case $\CA(G)$ fails by Theorem~\ref{cr14:thm:KMS}. That failure only removes
our obstruction; it does not construct a least representative. A possible
next step is to use that a putative omission set is computable from
\emph{every} computable-null erasure, a stronger demand than just $S\leT G$.

\begin{question}[Computational strength of null upper bounds]\label{cr14:q:cover}
Which degrees compute an $S\in\Znull$ such that $T\substar S$ for every
computable $T\in\Znull$? Does the answer change if $T$ is restricted to
one-point-per-block selectors for a fixed computable block layout?
\end{question}
Lemma~\ref{cr14:lem:cover} gives a necessary condition: such a degree computes
a computably bounded eventually different function, and hence is high or
DNR by Theorem~\ref{cr14:thm:KMS}. Sufficiency, the role of a fixed bound, and
the difference between all null sets and selectors should be investigated
separately. A proof for selectors alone would not automatically solve the
full null-cover question.

\begin{question}[How low can the starting object be?]\label{cr14:q:low}
Which nonzero degrees bound a binary set whose effective-dense class has no
least ordinary degree? Is there a c.e. binary counterexample? What effective
bounds can be imposed on descriptions witnessing its nontriviality?
\end{question}
Low $\Delta^0_2$ examples now exist by Theorem~\ref{cr14:thm:noleast}, so merely
asking for an arithmetical or halting-computable counterexample is no longer
a useful target here. The c.e. restriction is a different structural demand.

\subsection{Spectra and countable lower structure}
\begin{question}[Minimal representatives]\label{cr14:q:minimal}
Does either full class of a low $1$-generic contain a representative of
minimal ordinary Turing degree within that class?
\end{question}
The mask lattice and packed chain show extensive descent, but not descent
below an arbitrary representative. A promising restricted goal is to
characterize when a representative below $G$ has an effective-dense
equivalent strictly below it. Such a theorem would have to construct a new
representative below that prescribed oracle, rather than only one that
fails to compute it.

\begin{question}[Arbitrary countable coinitiality]\label{cr14:q:coinitial}
Can $[G]_{\ned}$ or $[G]_{\ued}$ have a countable coinitial family when $G$
is a low $1$-generic? Can the uniformity hypothesis of
Theorem~\ref{cr14:thm:countable} be removed or replaced by a weaker effective
presentation condition?
\end{question}
The packed chain has no lower bound in the class, but it need not be
coinitial. Theorem~\ref{cr14:thm:countable} eliminates uniformly supplied
descriptions; it does not select correct descriptions from an arbitrary
countable list of representatives. Pinning down that selection complexity
is a concrete intermediate target.

\begin{question}[The common lower cone of all computable erasures]
\label{cr14:q:common}
For a $1$-generic $G$, determine
\[
 \bigcap_{T\in\Comp\cap\Znull}\Low(G_T).
\]
When $G$ is low, is this ideal trivial? Can it be nonprincipal?
\end{question}
For two masks, Theorem~\ref{cr14:thm:masklattice} gives an exact answer.
For the full intersection, the present paper only excludes members that
compute an effective-dense description of $G$, under $\CA(G)$.
A common lower bound may carry information without itself providing any
such description.

\begin{question}[Classifying classes with least degrees]\label{cr14:q:positive}
Find intrinsic necessary and sufficient conditions for an effective-dense
class to have a least ordinary Turing degree, separately for the uniform
and nonuniform versions.
\end{question}
The criterion should distinguish genuine recoverability from every
\emph{description} from lower-boundedness among \emph{representatives}.
The repository's positive examples based on redundantly coded dyadic
columns provide a comparison family, but its coarse characterization
cannot simply be reused.\cite{ProveIt}

\subsection{Quantitative, effective, and formal directions}
\begin{question}[Effective presentations of descending chains]\label{cr14:q:chainindex}
What is the least possible complexity of a sequence of computable masks
whose erased oracles form a strict chain with no lower bound in the class?
Can a uniformly $G$-computable presentation occur without yielding a
uniformly $G$-computable sequence of descriptions?
\end{question}
Section~\ref{cr14:sec:chain} gives a $\doublezero$ presentation for the low
witnesses and carefully separates it from individual computability.
Determining the optimal bound requires an argument about presentations,
not just the degree of each separate mask.

\begin{question}[Other small-set ideals]\label{cr14:q:ideals}
For which effectively presented ideals $\mathcal I$ of negligible sets
can the sparse escape theorem and no-least argument be repeated, with
$\mathcal Z_0$ replaced by $\mathcal I$?
\end{question}
The proof suggests explicit hypotheses: computable masks should allow the
two description transformations; an $\mathcal I$-small omission set should
miss a point in every sufficiently late block of a chosen layout; and
computable selectors should lie in $\mathcal I$ with the desired budget.
Weighted-density and summable ideals are natural tests, but these hypotheses
must be checked rather than assumed.

\begin{question}[Randomness instead of genericity]\label{cr14:q:random}
Which randomness hypotheses give an analogue of
Lemma~\ref{cr14:lem:prediction} for computable coordinate erasures, and which
of those oracles also admit a suitable sparse escape theorem?
\end{question}
This separates a nonprediction problem from a highness problem.
Quasiminimality results for random sets in asymptotic degree structures do
not on their own settle this representative-spectrum question.\cite{AHJ}

\begin{question}[Binary representative complexity]\label{cr14:q:binary}
Can the description-selection machinery here contribute to a bound on the
complexity of a binary representative of a numerical function already
known to be nonuniformly effective-dense equivalent to some set?
\end{question}
This is a possible connection to Gerdes's Question~5 and repository C3,
not a solution of either. The present counterexamples start with a binary
set, so they avoid rather than solve that selection problem.\cite{Gerdes,ProveIt}

\begin{question}[Formalization of the proof interface]\label{cr14:q:lean}
Can generic nonprediction, budgeted sparse escape, and the class-level
leastness argument be formalized as reusable modules over a common oracle
computation interface in the repository?
\end{question}
A useful first milestone is a checked theorem taking $1$-genericity and
$\CA(G)$ as explicit hypotheses and proving Theorem~\ref{cr14:thm:noleast}.
The genericity constructions and the classical highness refinement should
then be proved or imported with their dependencies audited separately.

\section{Proof audit, formalization plan, and reproducibility}
\label{cr14:sec:audit}
\subsection{The dependency structure}
The shortest independent route to C2 has the following structure:
\[
 \begin{array}{c}
  2\text{-genericity}\Longrightarrow\CA(G)\\
  \CA(G)+\text{null omission set}\Longrightarrow
                   \text{computable sparse escape}\\
  1\text{-genericity}+\text{description below }G_T
                  \Longrightarrow T\substar S_d\\
  \text{least representative}\Longrightarrow\text{contradiction}.
 \end{array}
\]
A $2$-generic exists below $\doublezero$: enumerate the $\Sigma^0_2$ sets
of strings and, using $\doublezero$ to decide whether an extension exists,
meet the current set when possible and otherwise preserve an avoiding
prefix. Append a bit at each stage. This is the same meet-or-avoid
construction as in Proposition~\ref{cr14:prop:low}, one jump higher.
Hence the direct route has actual witnesses without appealing to
Theorem~\ref{cr14:thm:KMS}.

The low route instead uses
\[
 \text{low $1$-generic}\Longrightarrow
 \text{not high and not DNR}
 \overset{\text{classical Theorem~\ref{cr14:thm:KMS}}}{\Longrightarrow}\CA(G).
\]
The only named external classification in this route is stated in full
at its point of use. Neither route imports the repository's unrefereed
coarse-spectrum theorems as black boxes.

\subsection{A practical formalization boundary}\label{cr14:sec:formal}
An end-to-end formalization needs more than the quotient definition of
Turing degrees. The essential interfaces are finite-use oracle
computations with extension persistence; c.e. and $\Sigma^0_2$ predicates
on strings; explicit meet-or-avoid genericity; coding of total descriptions
with a distinguished omission symbol; and a density-zero library with
finite unions and prefix-count bounds.

The central theorem can first be proved for sets and total function
oracles, postponing passage to the degree quotient. Generic nonprediction
is a finite-computation argument. Budgeted blocks, selectors, and the
uniform P-ideal calculation are arithmetic constructions with eventual
bounds. The description-equivalence maps have direct pointwise proofs.
Only after these components are available does the class-level theorem
become a short composition of implications.

Theorem~\ref{cr14:thm:twoGenericCA} additionally needs the complexity calculation
for $B_\Gamma$ and the c.e. dense sets $U_k$. The low refinement needs
jump monotonicity, the generalized-low computation, and the classical
high-or-DNR characterization. Importing the latter as an axiom would make
the low result conditional in the formal development, even though it is
a published theorem in the mathematical argument. No Lean code in this
archive is represented as compiled or certified.

\subsection{What the accompanying finite checks do}\label{cr14:sec:checks}
The companion script \path{checks/check_finite.py} is deterministic and
uses only the Python standard library. It checks the two pointwise
correctness transformations for all small binary strings and masks;
finite erasure independence under a changed deleted bit; the join and
finite-amalgamation coordinate rules; capacity and freshness invariants of
the packed masks; and every-prefix budgets for representative finite block
layouts. It also checks finite instances of the dyadic-block inequality
used in Lemma~\ref{cr14:lem:Pideal}.

These checks are aimed at off-by-one, coordinate, and counting errors.
They do not test that a finite string is generic, exhibit a noncomputable
oracle in finite time, decide which programs are total, or verify an
infinite density limit. The proof of each such assertion is the
corresponding conventional argument above. The recorded output and a
machine-readable summary are supplied separately; rerunning the script
regenerates them.

The document is built by \path{build.sh}, which runs the checks and compiles
\path{article.tex} with pdfLaTeX. The archive contains the source, final
PDF, build instructions, source ledger, and proof audit. Fonts are supplied
by the user's TeX installation, not bundled in the archive.

\noindent\wtag\ In report~14 the script is \path{code/check_finite.py}; its
recorded outputs are \path{data/results.json} and \path{data/run.txt}
(byte-identical: the script prints the same text it writes). The script
writes \path{results.json} beside itself, and the delivered
\path{build.sh}, shipped as \path{code/build.sh}, changes to its own
directory, writes \path{checks/run.txt} there, compiles
\texttt{article.tex} and copies the PDF over \texttt{article.pdf}; neither
should be run in place. The delivered PDF is not shipped; the report's
\path{README.md} gives working commands on a copy.

\appendix
\section{Source provenance and novelty audit}
\label{cr14:app:sources}
\subsection{Repository snapshot}
The task was anchored to \repo\ commit
\begin{center}
\texttt{1a1396d4d3a2ac6812692df520517d86ba3a4785}.
\end{center}
The primary repository document is
\begin{quote}\small
\path{Computability/TuringDegrees/Research/CoarseDegrees/}\allowbreak
\path{research-plan/turing_degrees_unified.tex}.
\end{quote}
The fetched blob is
\texttt{ebecd2e15f8e4a06ebd0b0333387d27ebe630375}.
In that snapshot, source lines 552--573 include the established coarse
tools, the C2 question, its genericity sketch, and the explicit gap.
The C3--C8 follow-ups begin in the immediately following range.
The source describes itself as the second edition of a research synthesis;
it is not treated here as an externally refereed source.

\subsection{Primary literature checked}
Gerdes's arXiv record displayed version~1, submitted 9 August 2025, when
consulted on 28 September 2026. Section~7, Question~7 is the precise
least-representative question targeted here. The nearby questions about
binary representatives and descriptive complexity concern different
problems.\cite{Gerdes}

The Astor--Hirschfeldt--Jockusch preprint was checked for the definition of
effective-dense descriptions, the uniform and nonuniform reductions, and
the distinction between quasiminimality and internal ordinary-degree
structure. Its Theorem~5.9 is not cited as an internal least-degree theorem.
The paper's minimal-pair results also concern distinct questions about
asymptotic degrees.\cite{AHJ}

The Brendle--Brooke-Taylor--Ng--Nies paper records the equivalence between
computing an eventually different function and weak meager engulfing, and
attributes the exclusion of $2$-generics from that class to Rupprecht.
The direct proof in Section~\ref{cr14:sec:agreement} is included so that the
main application does not depend on translating that terminology.\cite{BBNN}

The high-or-DNR characterization was checked directly in
Kjos-Hanssen--Merkle--Stephan, Theorem~5.1, rather than inferred only from a
secondary reference. It is used solely for the non-high and low
refinements.\cite{KMS}

\subsection{What the audit does not establish}
Targeted searches for the combination of effective-dense representatives,
least ordinary Turing degrees, genericity, and eventually different
functions did not establish a prior occurrence of the exact erasure
application proved here. This is a limited search result, not evidence
that no earlier proof exists. In particular, elementary applications of
standard genericity may circulate in unpublished notes or use different
terminology. The mathematical claim is the stated theorem with its proof;
priority remains to be assessed independently.

\section{A compact checklist of delicate proof obligations}
\label{cr14:app:checklist}
\begin{longtable}{@{}>{\raggedright\arraybackslash}p{0.29\textwidth}p{0.66\textwidth}@{}}
\toprule
Issue & Where and how it is handled\\
\midrule
\endhead
Wrong values versus omissions & Lemma~\ref{cr14:lem:equivalence} gives two
explicit transformations. The inverse discards the whole computable mask.\\[4pt]
Full-class leastness & Theorem~\ref{cr14:thm:noleast} treats numerical
representatives, including those not below $G$, and uses binary comparison
witnesses in the uniform class.\\[4pt]
Nonuniform hard-coded data & The generic avoiding prefix, the computable
agreement index, and each finite list of total-function indices are
allowed to depend on the instance. No uniform selector is asserted.\\[4pt]
Decidable sparse mask & Equation~\eqref{cr14:eq:selector} gives one point in
each decidable interval using a total computable function, not a delayed
c.e. enumeration.\\[4pt]
Density at intermediate prefixes & Lemma~\ref{cr14:lem:blocks} treats every
$N\in[m_n,m_{n+1})$; Lemma~\ref{cr14:lem:Pideal} gives an explicit partial-block
bound.\\[4pt]
Bounded default in $f_S$ & Density zero makes the default branch in
\eqref{cr14:eq:choice} occur only finitely often. Infinite agreement remains
infinite after those exceptions are discarded.\\[4pt]
Infinite strict drops & The offsets $q_i(n)$ are fresh from all previously
inserted points; coincident selector functions cannot collapse a step.\\[4pt]
Countable presentations & Section~\ref{cr14:sec:countable} requires a uniform
$G$ procedure. Section~\ref{cr14:sec:chain} uses an explicitly nonuniform list
of total computable functions.\\[4pt]
Actual meet in the full degree order & The amalgamation argument of
Theorem~\ref{cr14:thm:masklattice} handles any common lower oracle, not just
another computable-mask oracle.\\[4pt]
Low refinement & Theorem~\ref{cr14:thm:KMS} is a named published input.
Generic non-DNR and generalized lowness are proved separately.\\[4pt]
No overclaim about minimality & The full-class theorem excludes least
elements; Question~\ref{cr14:q:minimal} keeps minimal elements separate.\\[4pt]
No experimental certification & Finite checks validate finite invariants
only. No independent review or proof-assistant verification is claimed.\\
\bottomrule
\end{longtable}

\begin{thebibliography}{9}\raggedright
\bibitem{ProveIt}
\repo\ repository,
\emph{Turing Degrees: A Unified Research Report}, prepared for Vladimir Reshetnikov,
research-plan source, second edition dated 18 September 2026,
repository snapshot \texttt{1a1396d4d3a2ac6812692df520517d86ba3a4785}.
\url{https://github.com/VladimirReshetnikov/ProveIt/blob/1a1396d4d3a2ac6812692df520517d86ba3a4785/Computability/TuringDegrees/Research/CoarseDegrees/research-plan/turing_degrees_unified.tex}.
Consulted 28 September 2026. Unrefereed research programme.

\bibitem{Gerdes}
Peter M. Gerdes,
\emph{Comparing Notions of Dense Computability on $\omega^\omega$ and
$2^\omega$}, arXiv:2508.06925v1, 2025.
\url{https://arxiv.org/abs/2508.06925v1}.
In particular, Section~7, Questions~5 and~7.

\bibitem{AHJ}
Eric P. Astor, Denis R. Hirschfeldt, and Carl G. Jockusch, Jr.,
\emph{Dense computability, upper cones, and minimal pairs},
Computability \textbf{8} (2019), 155--177.
DOI: \href{https://doi.org/10.3233/COM-180231}{10.3233/COM-180231}.
Author-hosted preprint:
\url{https://www.math.uchicago.edu/~drh/Papers/Papers/dense.pdf}.

\bibitem{KMS}
Bj\o rn Kjos-Hanssen, Wolfgang Merkle, and Frank Stephan,
\emph{Kolmogorov complexity and the Recursion Theorem},
Transactions of the American Mathematical Society \textbf{363} (2011),
5465--5480. Preprint arXiv:0901.3933v2.
\url{https://arxiv.org/abs/0901.3933v2}.
Theorem~5.1 is the high-or-DNR characterization used here.

\bibitem{BBNN}
J\"org Brendle, Andrew Brooke-Taylor, Keng Meng Ng, and Andr\'e Nies,
\emph{An analogy between cardinal characteristics and highness properties
of oracles}, arXiv:1404.2839v2, 2014.
\url{https://arxiv.org/abs/1404.2839v2}.
See Section~3.2, Theorem~4.2, and the discussion of $2$-genericity
in Section~4.
\end{thebibliography}
\end{document}
